\chapter{Workflows: Evaluate and Export}

\ChapterMeta{
This chapter outlines the canonical workflows for evaluating existing artifacts, exporting predictions from a backend, and executing convenient folder-level inference.
}{
It standardizes the generation of \path{predictions.json} and evaluation reports, ensuring that results remain strictly comparable and highly automatable.
}{
You will require a dataset root and one or more prediction artifacts; backend-specific operations may necessitate additional dependencies.
}
\ChapterAudience{Read this chapter when you already have data or model outputs and need to decide the fastest safe path to evaluation, export, or migration.}

\WorkflowOpening
{Pick the shortest route from existing outputs, backend inference, folder inference, or migration to a validated evaluation report.}
{Use this when you need to evaluate \path{predictions.json}, export one from a backend, inspect folder-level overlays, or normalize data from another ecosystem.}
{A dataset root or image folder; optionally an existing \path{predictions.json}, model artifact, backend runtime, or native dataset definition.}
{\cmd{yolozu validate predictions ... --strict} first, then the matching eval/export/migrate command for Workflow A--D.}
{\path{predictions.json}, validation output, evaluation JSON/HTML, overlays, or normalized dataset wrappers depending on the workflow.}
{Open the report JSON, confirm schema validation passed, and check that dataset split, class ids, and prediction counts match the intended run.}
{Missing files, stale relative paths, class-id drift, backend runtime imports, unsupported model shapes, or conceptual examples copied without real paths.}

\paragraph{Readiness classification.}
Workflow A in this chapter is the \textbf{stable main production lane} for \yolozu{}: validate predictions, then evaluate them under one pinned predictions interface contract.
Export paths and backend-specific execution remain useful, but they should be understood as secondary to that main lane.
See \path{docs/production_readiness.md}.

\paragraph{The shortest human summary.}
Most people arrive here with one of four questions:
\begin{itemize}
  \item ``I already have predictions. How do I evaluate them?''
  \item ``I have a model and a dataset. How do I export predictions safely?''
  \item ``I just want quick image-folder inference and overlays.''
  \item ``My data lives in another ecosystem. How do I migrate without rewriting everything?''
\end{itemize}

This chapter is organized exactly around those four questions.

\paragraph{Fast chooser.}
If you are unsure where to start:
\begin{itemize}
  \item choose \textbf{Workflow A} when inference already happened elsewhere,
  \item choose \textbf{Workflow B} when you need \yolozu{} to run inference and emit the standardized artifact,
  \item choose \textbf{Workflow C} when you want the quickest folder-level visual result,
  \item choose \textbf{Workflow D} when the problem is data or prediction migration rather than model execution.
\end{itemize}

\begin{figure}[h]
\centering
\resizebox{\linewidth}{!}{%
\begin{tikzpicture}[
  node distance=8mm,
  font=\small,
  box/.style={draw, rounded corners, align=left, inner sep=5pt},
  arrow/.style={-{Latex[length=2mm]}, thick}
]

% Row A
\node[box] (a_in) {\textbf{A}\\existing\\\path{predictions.json}};
\node[box, right=12mm of a_in] (a_val) {validate\\(\cmd{--strict})};
\node[box, right=12mm of a_val] (a_eval) {eval\\(suite / COCO)};
\node[box, right=12mm of a_eval] (a_rep) {reports\\JSON/HTML};
\draw[arrow] (a_in) -- (a_val);
\draw[arrow] (a_val) -- (a_eval);
\draw[arrow] (a_eval) -- (a_rep);

% Row B
\node[box, below=10mm of a_in] (b_in) {\textbf{B}\\backend\\inference};
\node[box, right=12mm of b_in] (b_exp) {export\\\path{predictions.json}\\+\cmd{export_settings}};
\node[box, right=12mm of b_exp] (b_val) {validate};
\node[box, right=12mm of b_val] (b_rep) {eval + reports};
\draw[arrow] (b_in) -- (b_exp);
\draw[arrow] (b_exp) -- (b_val);
\draw[arrow] (b_val) -- (b_rep);

% Row C
\node[box, below=10mm of b_in] (c_in) {\textbf{C}\\image folder\\or split};
\node[box, right=12mm of c_in] (c_run) {\cmd{predict-images}\\(+ overlays)};
\node[box, right=12mm of c_run] (c_out) {\path{predictions.json}\\+ HTML summary};
\draw[arrow] (c_in) -- (c_run);
\draw[arrow] (c_run) -- (c_out);

% Row D
\node[box, below=10mm of c_in] (d_in) {\textbf{D}\\dataset defs\\\path{data.yaml}/COCO};
\node[box, right=12mm of d_in] (d_mig) {migrate\\\path{dataset.json}\\labels/classes};
\node[box, right=12mm of d_mig] (d_out) {\path{predictions.json}\\(converted/exported)};
\node[box, right=12mm of d_out] (d_rep) {validate + eval};
\draw[arrow] (d_in) -- (d_mig);
\draw[arrow] (d_mig) -- (d_out);
\draw[arrow] (d_out) -- (d_rep);

\end{tikzpicture}
}
\caption{Workflow overview (A--D): keep inference flexible, fix the interface boundary at \cmd{predictions.json}.}
\end{figure}


\section{Workflow A: Evaluate existing predictions.json}
This is the safest and simplest path when a different system already produced model outputs.
You do not need to rerun inference; you only need a valid predictions artifact and a dataset root.

If you have already generated a \path{predictions.json} artifact from any inference backend, you can evaluate it directly:
\begin{lstlisting}[language=bash]
yolozu validate predictions /path/to/predictions.json --strict
python3 tools/eval_suite.py --dataset /path/to/yolo-dataset --predictions-glob '/path/to/predictions.json'
\end{lstlisting}

This represents the most common integration pathway for external inference systems.

\section{Workflow B: Run inference and export predictions}
Choose this path when you want \yolozu{} to be the boundary that turns backend-specific execution into a stable predictions artifact.

Depending on your deployment environment, several backends are available:
\begin{itemize}
  \item \textbf{Torch backend (research):} Exports predictions directly from a training checkpoint.
  \item \textbf{ONNX Runtime:} Facilitates CPU-friendly exports and rigorous parity checks.
  \item \textbf{OpenCV DNN (YOLO head):} Alternative ONNX execution via OpenCV DNN.
  \item \textbf{OpenCV DNN (RT-DETR, no NMS):} Logits+boxes decode path for RT-DETR ONNX (CUDA/CPU).
  \item \textbf{TensorRT:} Provides a pinned Linux/NVIDIA pipeline optimized for maximum throughput.
\end{itemize}

For production-oriented external inference, see \path{docs/external_inference.md} and the submodule-ready templates:
\path{examples/infer_cpp/} (CMake) and \path{examples/infer_rust/} (Rust).

\subsection{Rust template: optional ONNXRuntime mode}
The Rust template defaults to a minimal no-dependency interface-contract example, and can optionally enable ONNXRuntime via Cargo feature flags:
\begin{lstlisting}[language=bash]
# Minimal contract example (always available)
cargo build --release --manifest-path examples/infer_rust/Cargo.toml
examples/infer_rust/target/release/yolozu_infer_rust --out reports/pred_rust_stub.json

# Optional ONNXRuntime mode
cargo build --release --features onnxruntime --manifest-path examples/infer_rust/Cargo.toml
examples/infer_rust/target/release/yolozu_infer_rust \
  --mode onnxrt \
  --onnx /abs/path/model.onnx \
  --input-shape 1,3,64,64 \
  --image images/val/000001.jpg \
  --out reports/pred_rust_onnxrt.json
\end{lstlisting}

\noindent\textbf{Environment note:} The ONNXRuntime feature is intended for pinned Linux production images where Rust toolchains and ONNXRuntime shared libraries are provisioned explicitly.

\textbf{Conceptual (not copy-paste as-is):} repository-wrapper export shape.
Use real paths from \path{python3 -m yolozu export --help} for execution.
\begin{lstlisting}[language=bash]
# ONNXRuntime backend (input shape must match ONNX)
python3 -m yolozu export --backend onnxrt --dataset /path/to/yolo-dataset --split val --onnx /path/to/model.onnx --output reports/predictions.json

# TorchScript detection artifact with declared combined-output decode
python3 tools/export_predictions_torchscript.py --dataset /path/to/yolo-dataset --split val --model /path/to/model.torchscript --output reports/pred_torchscript.json --wrap

# ExecuTorch backend (.pte model; dry-run available for interface contract checks)
python3 -m yolozu export --backend executorch --dataset /path/to/yolo-dataset --split val --model /path/to/model.pte --output reports/pred_executorch.json --dry-run

# ExecuTorch non-dry decode: run ExecuTorch externally, then decode runtime rows
python3 -m yolozu export --backend executorch --dataset /path/to/yolo-dataset --split val --model /path/to/model.pte --runtime-output-json reports/executorch_runtime_outputs.json --boxes-scale norm --output reports/pred_executorch.json --force

# OpenCV DNN unified backend (auto decode + preset preprocess + IO dump)
python3 -m yolozu export --backend opencv-dnn --onnx /path/to/model.onnx --dataset /path/to/yolo-dataset --split val --max-images 8 --imgsz 640 --preprocess yolo_letterbox_640 --decode auto --dump-io reports/opencv_dump_io.json --output reports/pred_opencv_dnn_unified.json --force

# OpenCV DNN RT-DETR backend (no NMS)
python3 -m yolozu export --backend opencv-dnn-rtdetr --onnx /path/to/rtdetr.onnx --dataset /path/to/yolo-dataset --imgsz 640 --score-thr 0.01 --output reports/pred_rtdetr_opencv_backend.json --force

# OpenCV DNN YOLO backend (with NMS)
python3 -m yolozu export --backend opencv-dnn-yolo --onnx /path/to/yolo.onnx --dataset /path/to/yolo-dataset --imgsz 640 --score-thr 0.25 --nms-iou 0.45 --output reports/pred_opencv_dnn_yolo.json --force
\end{lstlisting}

For further details, consult:\\
\path{docs/training_inference_export.md}, \path{docs/external_inference.md}, and \path{docs/opencv_dnn_inference.md}.
For CI command-surface checks, prefer each exporter's \cmd{--dry-run} mode so schema and metadata contracts are validated without heavyweight runtime dependencies.

ExecuTorch non-dry mode intentionally does not pretend to execute a generic
mobile runtime inside YOLOZU. The declared path is: execute the \path{.pte}
artifact in the target ExecuTorch environment, write rows shaped
\cmd{[x1,y1,x2,y2,score,class_id]}, and pass that JSON via
\cmd{--runtime-output-json}. Unsupported row shapes fail before a predictions
artifact is written.

\subsection{Model intake before export (registry fetch)}
To standardize external checkpoints before running exporters, use the model fetch flow:
\begin{lstlisting}[language=bash]
yolozu list models
yolozu fetch yolox-s-coco --out models --accept-license
cat models/yolox-s-coco/meta.json
\end{lstlisting}

The generated \path{models/<id>/meta.json} always includes \path{source}, \path{version}, \path{license}, \path{sha256}, and \path{created_at}.

\subsection{YOLO user migration entrypoint (v5/v8/11/26)}
For users already operating within YOLO-family pipelines, \yolozu{} provides a direct interface-contract-first migration path:
\begin{lstlisting}[language=bash]
python3 tools/import_yolo_data_yaml.py --data-yaml /path/to/data.yaml --split val --output data/yolo_wrapper --force
python3 tools/export_predictions_yolo_runtime.py --model yolo11n.pt --dataset data/yolo_wrapper --split val --protocol nms_applied --wrap --output reports/pred_yolo_runtime.json
python3 tools/export_predictions_yolov5.py --dataset data/yolo_wrapper --split val --labels-dir runs/detect/exp/labels --protocol nms_applied --output reports/pred_yolov5.json
\end{lstlisting}

Protocol policy for fair comparisons:
\begin{itemize}
  \item \texttt{nms\_applied}: intended for YOLOv5/YOLOv8/YOLO11 post-NMS outputs.
  \item \texttt{e2e\_nms\_free}: intended for YOLO26/RT-DETR no-NMS outputs.
\end{itemize}

When outputs contain COCO \texttt{category\_id}, evaluate with class-id normalization:
\begin{lstlisting}[language=bash]
yolozu eval-coco --dataset data/yolo_wrapper --split val --predictions reports/pred_yolo_runtime.json --protocol nms_applied --classes data/yolo_wrapper/labels/val/classes.json --output reports/coco_eval_yolo_runtime.json
\end{lstlisting}

\subsection{YOLOX interoperability path}
For teams already operating YOLOX experiments (\path{exp.py} + checkpoints), keep native inference and export directly to the YOLOZU contract:
\begin{lstlisting}[language=bash]
python3 -m yolozu export --backend yolox --dataset /path/to/coco-yolo --split val2017 --exp /path/to/yolox_exp.py --weights /path/to/yolox_ckpt.pth --imgsz 640 --score-thr 0.01 --nms-iou 0.65 --output reports/pred_yolox.json
python3 tools/validate_predictions.py reports/pred_yolox.json --strict
python3 tools/eval_coco.py --dataset /path/to/coco-yolo --split val2017 --predictions reports/pred_yolox.json --protocol nms_applied --classes /path/to/coco-yolo/labels/val2017/classes.json --output reports/eval_yolox.json
\end{lstlisting}

Non-dry YOLOX export is fail-closed: \cmd{--exp} and \cmd{--weights} must
both identify existing files, the runtime must initialize, and inference must
run for every selected image. Otherwise the command returns nonzero without
writing a new predictions artifact. Use \cmd{--dry-run} only for a schema-valid
placeholder that intentionally skips inference.

The exporter records \path{meta.extra.execution_status},
\path{runtime_executed}, and \path{inference_calls} to distinguish
\cmd{dry_run} from \cmd{completed} execution. Exp/checkpoint paths and hashes
are stored under \path{meta.extra.model_provenance}; projected exp parameters,
preprocessing assumptions, and protocol metadata remain under
\path{meta.extra.export_settings}. A completed inference can validly contain
an empty detections list.

\subsection{YOLO-style external training lane}
When the goal is training rather than export-only interop, the recommended
external YOLO-style lane is YOLOX because it fits the repository's Apache-2.0
operational policy more naturally than optional copyleft-sensitive bridges.

\begin{lstlisting}[language=bash]
python3 -m yolozu train \
  --external-backend yolox \
  configs/examples/finetune_external/yolox_s_finetune_smoke.py \
  --dataset data/smoke \
  --split val \
  --dry-run \
  --output reports/train_external_yolox.json
\end{lstlisting}

The repo helper remains available when you want the lower-level bridge
surface directly:

\begin{lstlisting}[language=bash]
python3 tools/support_external_training.py train-yolox \
  --dataset data/smoke \
  --split val \
  --exp configs/examples/finetune_external/yolox_s_finetune_smoke.py \
  --dry-run \
  --output reports/support_external_training.train_yolox.json
\end{lstlisting}

\subsection{Current training support}
\begin{itemize}
  \item \textbf{Stable reference lane:} the in-repo RT-DETR pose reference trainer.
  \item \textbf{Supported external lane:} the Apache-2.0-friendly YOLOX path exposed through \cmd{yolozu train --external-backend yolox}.
  \item \textbf{Optional external bridges:} Ultralytics and HF DETR, kept explicit as external runtime/license boundaries.
  \item \textbf{Not claimed:} a universal training platform for every model family.
\end{itemize}

When post-train ONNX export is enabled, the export attempt runs on CPU even if training itself ran on CUDA or MPS. This keeps the export path more portable and avoids backend-specific exporter failures at the end of an otherwise successful run.

This writes a machine-readable report containing:
\begin{itemize}
  \item the resolved dataset root and split,
  \item a projected training configuration,
  \item the template command for the external launcher,
  \item the runtime/license boundary for the lane.
\end{itemize}

For real execution, provide your local YOLOX train launcher:
\begin{lstlisting}[language=bash]
python3 tools/support_external_training.py train-yolox \
  --dataset /path/to/coco-yolo \
  --split train \
  --exp /path/to/yolox_exp.py \
  --train-script /path/to/YOLOX/tools/train.py \
  --batch 16 \
  --weights /path/to/yolox_ckpt.pth \
  --output reports/support_external_training.train_yolox.exec.json
\end{lstlisting}

The Ultralytics bridge remains available on the same surface, but it is
documented as optional and should be reviewed separately under its own license
terms:
\begin{lstlisting}[language=bash]
python3 tools/support_external_training.py train-ultralytics \
  --dataset data/smoke \
  --split val \
  --preset smoke \
  --dry-run
\end{lstlisting}

The same top-level route also exposes the optional HF DETR bridge:
\begin{lstlisting}[language=bash]
python3 -m yolozu train \
  --external-backend hf-detr \
  facebook/detr-resnet-50 \
  --dataset data/smoke \
  --split val \
  --dry-run \
  --output reports/train_external_hf_detr.json
\end{lstlisting}

\subsection{SynthGen intake workflow (external generator)}
When synthetic shards are generated externally (for example in \path{YOLOZU-synthgen}), \yolozu{} provides the contract/eval side:
\begin{lstlisting}[language=bash]
python3 tools/validate_synthgen_contract.py \
  --input /path/to/synthgen_dataset/shards/train_000.jsonl \
  --max-samples 200

python3 tools/render_synthgen_overlay.py \
  --dataset-root /path/to/synthgen_dataset \
  --schema-id animal_v1 \
  --sample-index 0 \
  --output reports/synthgen_overlay_animal.png

python3 tools/eval_synthgen.py \
  --dataset-root /path/to/synthgen_dataset \
  --predictions /path/to/synthgen_dataset/shards/predictions_synthgen.json \
  --schema-id animal_v1 \
  --output reports/synthgen_eval_animal.json
\end{lstlisting}

Schema-specific task templates are provided under:
\begin{itemize}
  \item \path{configs/tasks/synthgen_animal_kpt.yaml}
  \item \path{configs/tasks/synthgen_mechanical_kpt.yaml}
\end{itemize}

For offline CI-safe smoke checks, use the bundled mini-shard fixture:
\begin{lstlisting}[language=bash]
python3 tools/smoke_synthgen.py \
  --dataset-root data/smoke/synthgen_minishard \
  --output-dir reports
\end{lstlisting}

Path-resolution note: if an external generator emits prediction rows with shard-relative asset paths (for example \path{../samples/...}), keep the predictions artifact under \path{shards/} or rewrite those paths before evaluation.

\section{Workflow C: Folder inference (pip CLI)}
Choose this when you care more about ``show me the result'' than about setting up a larger benchmark or protocol-pinned evaluation loop.

A highly convenient, user-facing approach for batch inference is:
\begin{lstlisting}[language=bash]
yolozu predict-images   --backend dummy   --input-dir /path/to/images   --output reports/predict_images.json   --overlays-dir reports/predict_images_overlays   --html reports/predict_images.html   --progress
\end{lstlisting}

This command generates a standardized predictions artifact, PNG visual overlays, and an HTML summary. With \cmd{--progress}, stderr shows stage and overlay progress. The terminal output prints the predictions path, HTML path, overlays directory, and first overlay path so new users know exactly what to open next. Use \path{configs/quickstart/predict_images_dummy.yaml} as the checklist for expected files.

\section{Post-result display and reporting}
Following inference or evaluation, YOLOZU generates both machine-readable artifacts and human-readable summaries.

\begin{longtable}{@{}p{0.34\textwidth}p{0.30\textwidth}p{0.30\textwidth}@{}}
\toprule
\textbf{Command family} & \textbf{Typical output} & \textbf{Usage} \\
\midrule
\cmd{yolozu predict-images} & \path{reports/predict_images.json} & Normalized predictions artifact \\
\cmd{yolozu predict-images} & \path{reports/predict_images.html} & Visual summary page \\
\cmd{yolozu predict-images} & \path{reports/predict_images_overlays/*.png} & Per-image visual overlays \\
\cmd{python3 tools/eval_suite.py} & \path{reports/eval_suite.json} & Protocol-pinned comparison report \\
\cmd{python3 tools/hpo_sweep.py} & \path{reports/hpo_sweep.jsonl} & Trial-level raw history \\
\cmd{python3 tools/hpo_sweep.py} & \path{reports/hpo_sweep.csv}, \path{reports/hpo_sweep.md} & Leaderboard and table views \\
\cmd{yolozu eval-instance-seg} & \path{reports/instance_seg_eval.json} & Instance segmentation metrics \\
\cmd{yolozu eval-instance-seg} & \path{reports/instance_seg_eval.html} & HTML report with optional overlays \\
\bottomrule
\end{longtable}

\subsection{Recommended result inspection sequence}
\begin{figure}[h]
\centering
\resizebox{0.95\linewidth}{!}{%
\begin{tikzpicture}[
  node distance=10mm,
  font=\small,
  box/.style={draw, rounded corners, align=left, inner sep=6pt},
  arrow/.style={-{Latex[length=2mm]}, thick}
]

\node[box] (v) {\textbf{1) Validate}\\schema + paths\\(\cmd{--strict})};
\node[box, right=16mm of v] (h) {\textbf{2) Review HTML}\\quick qualitative\\debug};
\node[box, right=16mm of h] (m) {\textbf{3) Inspect metrics}\\JSON/CSV/MD\\CI-friendly};

\draw[arrow] (v) -- (h);
\draw[arrow] (h) -- (m);

\end{tikzpicture}
}
\caption{Recommended result inspection order for fast iteration and reproducible reporting.}
\end{figure}

\begin{enumerate}
  \item \textbf{Validate the schema:}
  \begin{lstlisting}[language=bash]
yolozu validate predictions reports/predict_images.json --strict
  \end{lstlisting}
  \item \textbf{Review the HTML report} to perform rapid qualitative assessments.
  \item \textbf{Inspect the JSON/CSV/MD artifacts} for rigorous quantitative comparisons suitable for CI pipelines and academic publications.
\end{enumerate}

An example workflow for instance segmentation, incorporating HTML and overlay outputs:
\begin{lstlisting}[language=bash]
yolozu eval-instance-seg   --dataset examples/instance_seg_demo/dataset   --split val2017   --predictions examples/instance_seg_demo/predictions/instance_seg_predictions.json   --pred-root examples/instance_seg_demo/predictions   --classes examples/instance_seg_demo/classes.txt   --output reports/instance_seg_eval.json   --html reports/instance_seg_eval.html   --overlays-dir reports/instance_seg_overlays   --max-overlays 10
\end{lstlisting}

\section{Schema validation as a quality gate}
It is imperative to validate prediction artifacts prior to evaluation:
\begin{lstlisting}[language=bash]
python3 tools/validate_predictions.py /path/to/predictions.json --strict
\end{lstlisting}

This practice proactively prevents silent evaluation failures caused by malformed outputs.

\section{Workflow D: Migrate datasets and predictions}
Choose this path when the difficult part is not inference itself, but bridging from an existing dataset or prediction ecosystem into the \yolozu{} interface contract.

YOLOZU provides migration utilities that generate lightweight wrapper artifacts or converters, enabling seamless integration with prevalent ecosystem datasets. The core philosophy is to preserve the original dataset immutability while generating version-controlled wrappers.

\subsection{YOLO-style data.yaml / YOLO (YOLOv8/YOLO11)}
To generate a \path{dataset.json} wrapper from an existing \path{data.yaml}:
\begin{lstlisting}[language=bash]
yolozu migrate dataset   --from auto   --data /path/to/data.yaml   --output data/yolo_wrapper
\end{lstlisting}

Validate the resulting wrapper:
\begin{lstlisting}[language=bash]
yolozu validate dataset data/yolo_wrapper
\end{lstlisting}

\subsection{COCO JSON datasets (YOLOX / Detectron2 / MMDetection)}
To convert COCO instances JSON into YOLO-formatted label text files and generate a corresponding wrapper:
\begin{lstlisting}[language=bash]
yolozu migrate dataset   --from coco   --coco-root /path/to/coco_like_root   --split val2017   --output data/coco_yolo_like   --mode manifest
yolozu validate dataset data/coco_yolo_like --strict
\end{lstlisting}

To export that wrapper into external training layouts:
\begin{lstlisting}[language=bash]
yolozu export-dataset yolo   --dataset data/coco_yolo_like   --split val2017   --out-dir data/coco_yolo_export   --force
yolozu export-dataset coco   --dataset data/coco_yolo_like   --split val2017   --out-dir data/coco_export   --force
yolozu export-dataset kitti   --dataset data/coco_yolo_like   --split val2017   --out-dir data/coco_kitti_export   --force
\end{lstlisting}

For semantic-segmentation descriptors, export keeps the \path{dataset.json} contract while materializing \path{images/} and \path{masks/} into a lightweight external layout:
\begin{lstlisting}[language=bash]
yolozu doctor import --dataset-from auto --dataset /path/to/VOCdevkit/VOC2012 --split val --output -
yolozu import dataset   --from auto   --dataset /path/to/VOCdevkit/VOC2012   --split val   --output reports/voc_seg_wrapper   --force
yolozu export-dataset segmentation   --dataset reports/voc_seg_wrapper   --out-dir reports/voc_seg_export   --image-mode symlink   --force
\end{lstlisting}

For COCO keypoints roots that expose only \path{person_keypoints_<split>.json}, the same auto-detect preflight now recognizes the richer keypoints layout and writes a YOLOZU wrapper with \path{keypoint_names} / \path{skeleton} metadata before export:
\begin{lstlisting}[language=bash]
yolozu doctor import --dataset-from auto --dataset /path/to/coco_keypoints_root --split val2017 --output -
yolozu doctor train-dataset --from coco-keypoints --dataset /path/to/coco_keypoints_root --split val2017 --output -
yolozu import dataset   --from coco-keypoints   --dataset /path/to/coco_keypoints_root   --split val2017   --output reports/coco_keypoints_wrapper   --force
yolozu export-dataset coco   --dataset reports/coco_keypoints_wrapper   --split val2017   --out-dir reports/coco_keypoints_export   --force
\end{lstlisting}

For a repo-bundled tiny sample that exercises automatic detection plus \cmd{COCO -> YOLOZU -> YOLO/COCO/KITTI} end to end:
\begin{lstlisting}[language=bash]
yolozu validate dataset data/conversion_tiny_coco --split val2017 --strict
yolozu doctor import --dataset-from auto --dataset data/conversion_tiny_coco --split val2017 --output -
yolozu migrate dataset   --from auto   --dataset data/conversion_tiny_coco   --split val2017   --output reports/conversion_tiny_wrapper   --force
yolozu export-dataset yolo   --dataset reports/conversion_tiny_wrapper   --split val2017   --out-dir reports/conversion_tiny_yolo   --force
yolozu export-dataset coco   --dataset reports/conversion_tiny_wrapper   --split val2017   --out-dir reports/conversion_tiny_coco   --force
yolozu export-dataset kitti   --dataset reports/conversion_tiny_wrapper   --split val2017   --out-dir reports/conversion_tiny_kitti   --force
\end{lstlisting}

\subsection{Detectron2 / MMDetection predictions interop}
For Detectron2/MMDetection users, keep framework-native inference and only export detections into \path{predictions.json}:
\begin{lstlisting}[language=bash]
python3 tools/export_predictions_detectron2.py   --dataset data/coco_yolo_like   --split val2017   --config /path/to/d2_config.yaml   --weights /path/to/model_final.pth   --protocol nms_applied   --output reports/pred_detectron2.json

python3 tools/export_predictions_mmdet.py   --dataset data/coco_yolo_like   --split val2017   --config /path/to/mmdet_config.py   --checkpoint /path/to/epoch_12.pth   --protocol nms_applied   --output reports/pred_mmdet.json
\end{lstlisting}

Evaluate with explicit class-id normalization to avoid COCO \texttt{category\_id} drift:
\begin{lstlisting}[language=bash]
python3 tools/validate_predictions.py reports/pred_detectron2.json --strict
python3 tools/eval_coco.py   --dataset data/coco_yolo_like   --split val2017   --predictions reports/pred_detectron2.json   --protocol nms_applied   --classes data/coco_yolo_like/labels/val2017/classes.json   --output reports/coco_eval_detectron2.json
\end{lstlisting}

These commands are non-dry by default. Detectron2 requires existing
\cmd{--config} and \cmd{--weights} files; MMDetection requires existing
\cmd{--config} and \cmd{--checkpoint} files. Runtime initialization and
inference for every selected image must complete before a new artifact is
written. Missing prerequisites, unreadable inputs, or runtime/inference errors
return nonzero. Use \cmd{--dry-run} only for schema-valid output that
intentionally skips framework inference.

The exporters record \path{meta.extra.execution_status},
\path{runtime_executed}, and \path{inference_calls}; successful non-dry
artifacts also include config/checkpoint provenance and SHA-256 values in
\path{meta.extra.model_provenance}. Empty detections remain valid after
completed inference. The exporters preserve raw framework box convention
(\cmd{xyxy_abs}) in \path{export_settings.raw_output_bbox_format}, while
normalizing the predictions interface contract output to \cmd{cxcywh_norm} for
strict schema validation.

\subsection{COCO results (inference outputs) \texorpdfstring{$\rightarrow$}{->} predictions.json}
To convert COCO-style detection results into the standardized YOLOZU predictions contract:
\begin{lstlisting}[language=bash]
yolozu migrate predictions   --from coco-results   --results /path/to/coco_results.json   --instances /path/to/instances_val2017.json   --output reports/predictions.json   --force

yolozu validate predictions reports/predictions.json --strict
\end{lstlisting}
